% introduction.tex -- Cross-Enterprise Verification Paper (RU-V7)

\section{Introduction}

Enterprise blockchain networks increasingly serve multi-party ecosystems where
independent organizations---manufacturers, suppliers, logistics providers,
financial institutions---must verify each other's state transitions without
exposing proprietary operational data. In a supply chain setting, a
manufacturer must prove that goods recorded in its internal ledger correspond
to a purchase order in the supplier's ledger, while neither party reveals
inventory levels, pricing, or transaction volumes to the other or to the
settlement layer. Zero-knowledge validium architectures, where each enterprise
maintains private off-chain state and submits succinct validity proofs to a
shared L1~\cite{ethereum2024validium}, provide the foundation for such
privacy-preserving verification. However, existing validium systems treat each
enterprise as an isolated rollup: there is no standard mechanism for verifying
\emph{cross-enterprise} interactions---transactions that span two or more
private state trees---without breaking the privacy boundary.

The problem is not merely theoretical. Production multi-chain environments
such as Polygon AggLayer~\cite{polygon2025agglayer} and zkSync
Gateway~\cite{zksync2025gateway} aggregate proofs from multiple chains to
reduce L1 verification cost, but they assume public or semi-public state and
do not address the case where each chain's data must remain confidential.
Rayls~\cite{rayls2025privacy}, developed by Parfin for the Drex
(Brazilian CBDC) ecosystem, introduces private subnets with ZK proofs and
homomorphic encryption, but relies on a trusted hub chain that sees
transaction metadata. zkCross~\cite{zkcross2024} proposes cross-chain
privacy-preserving auditing but targets compliance rather than real-time
operational verification.

This paper presents a \emph{cross-enterprise verification protocol} for ZK
validium systems that enables two or more enterprises to prove the consistency
of shared interactions---a sale recorded in both parties' state trees, a
cross-reference between supply chain events---without revealing any data
beyond the mere existence of the interaction. The protocol introduces a
\emph{cross-reference circuit} that verifies Merkle inclusion proofs from two
independent sparse Merkle trees (SMTs) and binds them through a Poseidon
commitment, using only already-public state roots as public inputs. The
information leakage is exactly 1 bit per interaction: the fact that an
interaction exists. No keys, values, amounts, or transaction details are
exposed.

We evaluate three verification approaches---Sequential, Batched Pairing, and
Hub Aggregation---across enterprise counts ranging from 2 to 50, with
interaction densities from sparse (linear chain) to dense (interactions
exceeding enterprise count). Our key finding is that cross-enterprise
verification achieves $1.41\times$ overhead in the sequential case and
$0.64\times$ overhead with batched pairing verification, both well below the
$2\times$ target. The cross-reference circuit requires 68,868 constraints
(two depth-32 Merkle path verifications plus a Poseidon-4 commitment) and
generates proofs in 448~ms with rapidsnark. These results demonstrate that
privacy-preserving cross-enterprise verification is practical for
enterprise validium deployments with sub-second proof generation and
manageable L1 gas overhead.
